\subsection{\problemblock{*Limits} Data Block}\label{sec:block-limits}

TAOS, being a three-degree-of-freedom simulation, assumes that a control system is
available that can achieve the body attitudes requested in the guidance scheme.
Usually the control system can only maintain body attitudes for certain flight
conditions, and the body-attitude angles are restricted as well. The
\problemblock{*limits} data block provides a way to keep the trajectory within
these bounds.

Limits can be set on any of the guidance variables given in
\cref{tab:body-attitude-guidance-angles,tab:flight-condition-guidance}, except for
\texttt{intercept}, \texttt{prop}, \texttt{downria}, \texttt{upria}, and
\texttt{l/d\_max}. For example,

\begin{lstlisting}[style=taosmanual]
*limits alpha<20 alpha>-10
\end{lstlisting}

puts a $20^{\circ}$ upper boundary and a $-10^{\circ}$ lower boundary on angle of
attack. Greater-than and less-than signs are used to designate whether the limit is an
upper or lower bound.

The limiting variables do not have to be the same as those used for the guidance
scheme; limits can be imposed on any of the guidance variables. Any number of
limits can be specified. For example,

\begin{lstlisting}[style=taosmanual]
*limits bankgd>-60 bankgd<60 alpha>-15 alpha<15.0
        beta>-15 beta<15
\end{lstlisting}

puts limits on three of the control variables: \statevar{alpha}, \statevar{beta}, and
\statevar{bankgd}.

Flight-condition guidance (\cref{tab:flight-condition-guidance}) requires iterative
searches to solve for the control variables. Limits on the body-attitude angles help
control convergence during these searches.

\begin{quote}\small\itshape
Source terminology. The source page prints \texttt{prop} in the list of guidance rules that cannot be
limited. The corresponding guidance-rule name elsewhere in the manual is
\texttt{propnav}; the source wording is retained here and the relationship is recorded
in the Chapter~4 editorial metadata.
\end{quote}
